% Einheit 3 von 3 – Unabhängigkeit als Argument: die Ausgleichs-Fehlvorstellung widerlegen (bank/unabhaengigkeit/e3.jsonl, zusammenbau v0.1)
%% TODO Zeitmarke Einheit 3: Marken-Zeile nicht lesbar
%% TODO Zweigzeile Teil 1: was hier gelernt wird (Satz in Schülersprache aus dem Einheitstitel) – Einheit 3
\einheitenkopf[e3]{Einheit 3 von 3 · Unabhängigkeit als Argument: die Ausgleichs-Fehlvorstellung widerlegen}
\zweigzeile{Abitur GK}
\setcounter{aufgabe}{13}

%% TODO Titel: Ich-kann-Satz fehlt in der Bank (Platzhalter: Kettenname) – Nr. 14
%% TODO Anweisung: ein Satz über der ersten Teilaufgabe fehlt in der Bank – Nr. 14
\begin{aufgabe}{Als Argument}
%% TODO \rechenplatz{n}: Schrittzahl des Grundfalls fehlt in der Bank (nur bei fester Schreibform, 2.3 b)
\begin{teile}
\teil Eine Münze zeigt dreimal hintereinander Wappen. Wie wahrscheinlich ist beim nächsten Wurf Zahl? \\ \kreuz{$0{,}5$} \\ \kreuz{größer als $0{,}5$} \\ \kreuz{kleiner als $0{,}5$}
\teil Ein Glücksrad zeigt bei jeder Drehung mit $25\,\%$ Wahrscheinlichkeit Blau. Bei 80 Drehungen erschien Blau nur 12-mal. Mia folgert, dass bei den nächsten 80 Drehungen der Anteil von Blau deutlich über $25\,\%$ liegen muss. Beurteile die Aussage \hfill (Abitur 2017 LK).
\teil Bei einem Online-Quiz beantwortet Jan jeden Samstag 8 Fragen; jede Frage beantwortet er unabhängig von den anderen mit der Wahrscheinlichkeit $0{,}6$ richtig. Die Anzahl der richtigen Antworten an einem Samstag ist binomialverteilt. An drei Samstagen hintereinander hatte er jeweils 7 richtige. Jan meint, die Chance auf 7 richtige sei an diesem Samstag deshalb deutlich kleiner. Beurteile die Aussage \hfill (Abitur 2021 GK).
\end{teile}
\end{aufgabe}

%% TODO Titel: Ich-kann-Satz fehlt in der Bank (Platzhalter: Kettenname) – Nr. 15
%% TODO Anweisung: ein Satz über der ersten Teilaufgabe fehlt in der Bank – Nr. 15
\begin{aufgabe}{Als Argument – Fehler finden, Begründen}
\begin{teile}
\teil Nach 50 Münzwürfen mit nur 18-mal Wappen sagt Paul: „Das Gesetz der großen Zahlen sagt, dass sich die Häufigkeit bei der Hälfte einpendelt – also muss jetzt öfter Wappen kommen.“ Finde den Fehler in der Begründung.
\teil Begründe, warum die Wahrscheinlichkeit für Wappen beim nächsten Münzwurf nicht davon abhängt, was in den bisherigen Würfen gefallen ist.
\end{teile}
\end{aufgabe}
